\documentclass[11pt]{article}
\usepackage{graphicx} % Required for inserting images
\usepackage{hyperref} 
\usepackage{fullpage}
\usepackage{comment}
\usepackage[]{natbib}
\hypersetup{
    colorlinks=true,
    linkcolor=blue,  
    citecolor=black,
    urlcolor=blue        
}



\begin{document}

\title{Supplemental Resources}
\maketitle

Supplementary resources related to this article are accessible at DOI 10.17605/OSF.IO/JQNHG \cite{LaCombe_Lavin_Horton_Ball_Holler_2025}

The following annotated bibliography is intended to provide an overview of key papers for further reading into particular R\&R principles and examples mentioned in this paper.

\vspace{.1in}
\noindent
{\bf Active Learning:}

Shi, Y., Yang, H., MacLeod, J., Zhang, J., \& Yang, H. H. (2020). College Students’ Cognitive Learning Outcomes in Technology-Enabled Active Learning Environments: A Meta-Analysis of the Empirical Literature. \emph{Journal of Educational Computing Research}, 58(4), 791–817. \url{https://doi.org/10.1177/0735633119881477}

Tutal, {\"O}., \& Yazar, T. (2022). Active learning promotes more positive attitudes towards the course: A meta‐analysis. \emph{Review of Education}, 10(1), e3346. \url{https://doi.org/10.1002/rev3.3346}

\vspace{.1in}
\noindent
{\bf Peer Instruction:} 

Crouch, C. H., \& Mazur, E. (2001). Peer Instruction: Ten years of experience and results. \emph{American Journal of Physics}, 69(9), 970–977. \url{https://doi.org/10.1119/1.1374249}

\vspace{.1in}
\noindent
{\bf Meta-cognition:}

Fleur, D. S., Bredeweg, B., \& van den Bos, W. (2021). Metacognition: Ideas and insights from neuro- and educational sciences. \emph{Npj Science of Learning}, 6(1), 1–11. \url{https://doi.org/10.1038/s41539-021-00089-5}

Taczak, K., \& Roberston, L. (2017). Chapter 11. Metacognition and the Reflective Writing Practitioner: An Integrated Knowledge Approach. In P. Portanova, J. M. Rifenburg, \& D. Roen (Eds.), Contemporary Perspectives on Cognition and Writing (pp. 211–229). The WAC Clearinghouse; University Press of Colorado. \url{https://doi.org/10.37514/PER-B.2017.0032.2.11}

\vspace{.1in}
\noindent
{\bf Effective Peer Review and Self-Assessment:}

Gorzelsky, G., Driscoll, D. L., Paszek, J., Jones, E., \& Hayes, C. (2016). Chapter 8. Cultivating Constructive Metacognition: A New Taxonomy for Writing Studies. In C. Anson \& J. L. Moore (Eds.), Critical Transitions: Writing and the Question of Transfer (pp. 215–246). The WAC Clearinghouse; University Press of Colorado. \url{https://doi.org/10.37514/PER-B.2016.0797.2.08}

\vspace{.1in}
\noindent
{\bf Reproducibility \& Replicability:}

The Turing Way, \url{https://book.the-turing-way.org}

A group of psychologists (https://osf.io/eg4hp ) suggested these resources:

Lindsay et al. (2016), Research Preregistration 101

Galetzka, C. (2019). A Short Introduction to the Reproducibility Debate in Psychology. \emph{Journal of European Psychology Students}, 10(1), 16–25. \url{https://doi.org/10.5334/eps.469}

Aczel et al. (2019). A Consensus-based Transparency Checklist. \emph{Nature Human Behavior},
4, 4–6, \url{https://doi.org/10.1038/s41562-019-0772-6}

Soderberg, C. K. (2018). Using OSF to Share Data: A Step-by-step Guide.
\emph{Advances in Methods and Practices in Psychological Science}, 1 (1), 115–120, \url{https://doi.org/10.1177/2515245918757689}


\vspace{.1in}
\noindent
{\bf FAIR Principles:}

GO FAIR Initiative, \href{https://www.go-fair.org}{https://www.go-fair.org} 

\vspace{.1in}
\noindent
{\bf Code Smells/Code Quality/Code Review:}

Fowler, M., \& Beck, K. (1999). Refactoring: Improving the Design of Existing Code. Addison-Wesley Professional.

Thomas, D. \& Hunt, A. (2019). The Pragmatic Programmer: Your Journey to Mastery, 20th Anniversary Edition. Addison-Wesley Professional.

Pruim, Randall, Gîrjău. M.-C. \& Horton, N. J. (2023), ‘Fostering Better Coding Practices for Data Scientists’, \url{https://doi.org/10.1162/99608f92.97c9f60f}.

\vspace{.1in}
\noindent
{\bf Literate Programming:}

Kery, M. B., Radensky, M., Arya, M., John, B. E., \& Myers, B. A. (2018, April 19). The Story in the Notebook: Exploratory Data Science using a Literate Programming Tool. Proceedings of the 2018 CHI Conference on Human Factors in Computing Systems. CHI ’18: CHI Conference on Human Factors in Computing Systems, Montreal QC Canada. \url{https://doi.org/10.1145/3173574.3173748}

Knuth, D. E. (1984). Literate programming. \emph{The Computer Journal}, 27(2), 97–111.

Baumer, B., Cetinkaya-Rundel, M., Bray, A., Loi, L., and Horton, N.J. (2014). R Markdown: Integrating A Reproducible Analysis Tool into Introductory Statistics.
Technology Innovations in Statistics Education, \url{https://escholarship.org/uc/item/90b2f5xh}.

\vspace{.1in}
\noindent
{\bf Portfolio Model:}

Arola, K. L., Ball, C. E., \& Sheppard, J. (2014). Writer/Designer: A Guide to Making Multimodal Projects. Macmillan Higher Education.

\vspace{.1in}
\noindent
{\bf Reproducibility Audits:}

Nüst, D., Ostermann, F. O., Sileryte, R., Hofer, B., Granell, C., Teperek, M., Graser, A., Broman, K., Hettne, K., Clare, C. Belliard, F., Wang, Y. (2024). AGILE Reproducible Paper Guidelines. https://doi.org/10.17605/OSF.IO/CB7Z8.

Le, Benjamin (2025). PSYC 226 - Open Science \& Inclusive Psychology: Assignments. OSF. https://doi.org/10.17605/OSF.IO/4GYQX 

Berkeley Initiative for Transparency in the Social Sciences (2020). Guide for Advancing Computational Reproducibility in the Social Sciences. https://bitss.github.io/ACRE/.

\vspace{.1in}
\noindent
{\bf Scaffolding:}

Saeed, M., \& Fund, F. (2024). Scaffolding Reproducibility for the Machine Learning Classroom. ACM. \url{ttps://acm-rep.github.io/2024/posters/ACM_REP24_paper_29.pdf}


\bibliography{bibliography}{}
\bibliographystyle{agsm}

\end{document}
